\section{Results}
\label{sec:results}

\subsection{Surface-water engine performance}
Before the groundwater leg is added, the surface-water PFAS engine is validated on its own across
the full Rogue network. The implementation passes its correctness and conservation checks: the soil
three-phase equilibrium solver reproduces a 40-digit decimal reference to a worst-case relative error
of $1.7\times10^{-7}$ over $29{,}160$ parameter combinations, the in-stream routine agrees with the
established SWAT\textsuperscript{+} pesticide channel solver (reduced to its linear-partition subset)
to a maximum absolute difference of $2.4\times10^{-7}$~kg, and the land- and channel-phase mass
balances close to a relative residual of order $10^{-5}$. The uncalibrated model, initialized from
land-use and soil-class background concentrations alone, reproduces the spatial \emph{pattern} of
ambient PFOS --- including the mainstem downstream gradient --- but overpredicts the \emph{magnitude}
by a nearly uniform factor of approximately nine. Because reach concentration is linear in soil
loading, a single global soil-loading scale factor of $0.11$ corrects the magnitude in one step,
reducing the log-space root-mean-square error across the $20$ gauged reaches ($29$ stations) from
$1.33$ to $0.74$~dex and the percent bias from $+1516\%$ to $+78\%$. The improvement is concentrated
where the model is physically anchored: along the seven source-bearing mainstem reaches the
calibrated log-RMSE is $0.15$~dex (a factor of $1.4$) with a log-space Nash--Sutcliffe efficiency of
$+0.61$ and a percent bias of $-7\%$, and the modeled concentration rises monotonically downstream
from $6.2$~ng~L\textsuperscript{-1} in the headwaters toward the Wolverine corridor against an
observed rise from $6.3$ to $27$~ng~L\textsuperscript{-1}. The calibrated in-stream PFOS field
renders this gradient over the full $590$-reach network (Figure~\ref{fig:instream}), and the
remaining misfit is structured rather than random: the two lowest mainstem reaches underpredict the
observed peak ($14.7$--$15.4$ vs $19$--$27$~ng~L\textsuperscript{-1}), the signature of the
concentrated House Street source that a uniform soil background cannot reproduce, while the
bidirectional tributary scatter reflects per-land-use heterogeneity in soil PFAS. These per-HRU
inputs --- the initial soil PFOS pool and the discrete, soil-class-controlled Freundlich sorption
fields the engine reads (Figure~\ref{fig:soil}) --- are the national, automatable initialization
that replaces the measured per-site soil inventory available only for data-rich legacy basins.

\begin{figure}[ht]
  \centering
  \includegraphics[width=\textwidth]{fig2_instream_pfos.pdf}
  \caption{Calibrated ($0.11$ soil-loading scale, no point source) flow-weighted in-stream PFOS over
  the $590$-reach Rogue network on a logarithmic colour scale. Concentration rises along the mainstem
  from green tributaries through the orange mid-river toward the Wolverine source (red star); open
  rings are the EGLE ambient monitoring stations.}
  \label{fig:instream}
\end{figure}

\begin{figure}[ht]
  \centering
  \includegraphics[width=\textwidth]{fig4_soil_params.pdf}
  \caption{Per-HRU soil-PFAS initialization across the $17{,}771$ Rogue HRUs: (a) initial
  surface-layer soil PFOS pool (continuous, NaturalBreaks classes); (b) Freundlich sorption
  coefficient $k_f$ and (c) Freundlich exponent $n$, both discrete soil-class fields. Colour map:
  batlow (perceptually uniform, colour-blind safe). The three panels share the same spatial extent;
  the scale bar and north arrow on panel~(a) apply to all.}
  \label{fig:soil}
\end{figure}

\subsection{Surface-water PFAS}
The surface-water engine reproduces ambient PFOS on the Rogue from the single global soil-loading
scale factor calibrated above: along the seven source-bearing mainstem reaches the calibrated
log-space root-mean-square error is $0.15$~dex (a factor of $1.4$), and over all $20$ gauged reaches
it is $0.74$~dex (Table~\ref{tab:validation}). The model captures the mainstem downstream gradient
toward the Wolverine corridor (rising from $\sim$6 to $\sim$27~ng~L\textsuperscript{-1}); the
all-network residuals are structured rather than random, localizing the concentrated point source and
per-land-use soil heterogeneity. This land-derived in-stream PFAS load is this paper's
surface-water-alone result and is the baseline to which the groundwater pathway is added below
(Figure~\ref{fig:swpfas}).

\begin{figure}[ht]
  \centering
  \includegraphics[width=\textwidth]{fig3_calibration.pdf}
  \caption{Surface-water PFAS on the Rogue. (a) Observed versus modeled in-stream PFOS, one point
  per gauged reach; the source-bearing mainstem reaches lie within the $\times2$ envelope while
  tributaries scatter (the structured residual that localizes source and soil heterogeneity).
  (b) The ambient monitoring stations coloured by observed PFOS, showing the downstream rise toward
  the Wolverine source.}
  \label{fig:swpfas}
\end{figure}

\subsection{Groundwater model generation and flow calibration}
The automated pipeline generates a complete, calibrated MODFLOW~6 groundwater flow model for the
Rogue with no manual editing, yielding a simultaneously head- and flux-consistent field: the
pilot-point inversion with a baseflow constraint reaches a head Nash--Sutcliffe efficiency of
$\mathbf{0.91}$ (RMSE $5.6$~m over $5{,}383$ static water-table heads) while reproducing the observed
gaining baseflow ($+5.46$ vs $5.56$~m\textsuperscript{3}~s\textsuperscript{-1};
Table~\ref{tab:validation}). The model comprises $\approx 24{,}900$ active cells and a connected SFR
network of $1{,}506$ reaches built directly from the SWAT\textsuperscript{+} channel graph.
Reaching this on a second, heavily-developed and dissected watershed --- beyond the
lightly-developed basin on which the generator was first built --- requires three reproducible
ingredients: scaling raw well pump \emph{capacity} to actual withdrawal, a near-no-flow
watershed-divide boundary, and a uniformly-thick bedrock layer for steady-state Newton stability.
Two calibrated multipliers carry the weight of these adjustments and should be read as effective, not
literal, values (Table~\ref{tab:calparams}): the well-pumping multiplier settles near $-1.26$ in
$\log_{10}$ --- actual withdrawal $\sim$5\% of raw registered capacity, absorbing the gap between
permitted capacity and long-term-average use --- and the drain conductance rails at its upper prior
bound, standing in for unresolved near-surface drainage structure rather than a measured conductance.

\subsection{Groundwater PFAS plume, validated against the Wolverine measurements}
The Freundlich GWT model reproduces the House Street plume on the calibrated flow field, validated
against $846$ measured groundwater PFOS observations. The measurements aggregate to $73$ grid cells
(the per-cell median; ``$846$ observations'' is the data volume, not the independent degrees of
freedom), of which $8$ source-zone cells are \emph{prescribed} and the remaining $65$ are
\emph{predicted}. Over all $65$ predicted cells the log-space RMSE is $1.1$~dex with $71\%$ within a
factor of ten, and among the $26$ predicted cells \emph{above} background (observed
$>50$~ng~L\textsuperscript{-1}) it is $1.4$~dex with $46\%$ ($12$ of $26$) within a factor of ten
(Table~\ref{tab:validation}; Figure~\ref{fig:gwpfas}; Supplementary Figure~S5). The source is
\emph{prescribed} to the measured plume (constant-concentration cells set to observed values
$>10^4$~ng~L\textsuperscript{-1}), so the skill assessment separates prescribed from predicted cells:
the model captures the plume \emph{shape near the source} and predicts ambient background away from
it. The all-cell figure is flattered by the many background-dominated cells where both model and
observations sit near the ambient floor; the above-background figure is the more demanding test. The
elevated observations distant from House Street indicate additional source areas the single
prescribed source does not represent, not a transport error. At this resolution a grid P\'eclet
number of $\sim$25 ($\Delta x/\alpha_L = 250/10$) means numerical dispersion dominates the prescribed
physical dispersivity. The streamflow-transport package routes the groundwater-discharged PFAS
downstream through the channel network.

A more direct test validates the pathway at the discharge interface itself. Streambed pore water along
the lower mainstem ($43$ sites, all within $100$~m of a channel) measures the groundwater
concentration where it enters the stream, and its skill is sharply concentration-dependent: below
$10$~ng~L\textsuperscript{-1} the model is essentially unbiased ($-0.19$~dex, $82\%$ within a factor of
ten), whereas above $50$~ng~L\textsuperscript{-1} it under-predicts systematically ($-1.43$~dex). The
deficit is therefore not random and not a near-detection-limit artefact: it is localized to the
highest-concentration discharge reaches near the source, where the $250$~m cell smears the steep plume
gradient --- the same numerical-dispersion limit identified above, now pinned to the reaches that carry
the bulk of the discharged mass. Reporting a single pooled error ($1.2$~dex) would hide this structure
and is avoided here.

Because the comparison above is in concentration, it understates the mass implication. The conserved
quantity is load (concentration times the gaining flux): the net groundwater-to-stream exchange
integrates to $+5.46$~m\textsuperscript{3}~s\textsuperscript{-1}, matching the calibrated baseflow, and
the transport-delivered groundwater PFOS load to the network is of order $45$~kg~yr\textsuperscript{-1}.
Two reaches of equal concentration but differing gaining flux carry proportionally different mass, so
the validation and the surface/groundwater apportionment are assessed in load space, not concentration
alone.

\begin{figure}[ht]
  \centering
  \includegraphics[width=\textwidth]{rogue_pfas_validation.png}
  \caption{Groundwater PFAS on the Rogue. (a) Modeled groundwater PFOS plume with the $846$
  measured observations overlaid (same colour scale) and the House Street source; (b) modeled
  versus observed groundwater PFOS, one point per validation cell, with the $1{:}1$ line and
  $\times10$ envelope (panel-title log-RMSE $\approx1.1$~dex over all predicted cells, rounded; the
  honest above-background figure is $1.4$~dex, Section~\ref{sec:pfas} and Supplementary Figure~S5);
  (c) the groundwater-discharged PFOS routed by SFT into $1{,}354$ of the $1{,}506$ stream reaches
  (raw, before dilution by surface streamflow).}
  \label{fig:gwpfas}
\end{figure}

\subsection{Detecting the groundwater contribution model-free: multi-analyte source fingerprinting}
Before fitting any model to the in-stream data, we test whether the groundwater contribution is
detectable in the observations themselves. PFAS compounds of different perfluorocarbon chain length
partition differently between the surface and subsurface pathways --- short-chain compounds are weakly
sorbed and mobile, long-chain compounds are strongly retained --- so the legacy groundwater plume and
the diffuse surface signal carry distinct chemical fingerprints. Across the watershed the two
end-members are sharply separated (Figure~\ref{fig:finger}a): the groundwater plume is PFOS-dominated
(PFOS is $0.76$ of the five-analyte sum of PFOS, PFOA, PFHxS, PFBS and PFHxA), whereas the in-stream
surface signal is short-chain-rich (PFBS and PFHxA together $0.42$, versus $0.05$ in the plume). Their
composition distance is near-maximal ($L_1=0.98$ of a possible $2$), so an in-stream sample is an
identifiable two-end-member mixture.

Decomposing each mainstem in-stream sample into the surface and groundwater end-members by non-negative
unmixing --- introducing no fitted transport parameter --- establishes a clear downstream gradient
(Figure~\ref{fig:finger}b): the groundwater fraction rises from $\approx0$ in the headwaters to
$0.55$--$0.63$ at the lower mainstem near the Wolverine source. This decomposition is model-free ---
unlike a model fit it does not rest on the aquifer-side validation, and unlike a single-analyte
comparison it cannot be reproduced by any co-located surface source with a different chemical signature.
It is thus the primary evidence that the lower-mainstem PFOS is groundwater-borne; the joint calibration
below then \emph{quantifies} that contribution and independently recovers the same spatial gradient
(Pearson $r=0.58$).

\begin{figure}[ht]
  \centering
  \includegraphics[width=\textwidth]{fig_fingerprint_separation.png}\\[2pt]
  \includegraphics[width=0.7\textwidth]{fig_reach_gw_fraction_vs_paper.png}
  \caption{Multi-analyte source fingerprinting. (a) Mean five-analyte composition of the groundwater
  plume, ambient groundwater, and in-stream surface water; the plume is PFOS-dominated and the surface
  signal short-chain-rich. (b) Two-end-member mixing space (PFOS fraction versus short-chain fraction);
  in-stream samples lie on the surface--plume mixing line. (c) Groundwater fraction of in-stream PFAS
  along the Rogue mainstem from the parameter-free fingerprint decomposition (green) versus the
  joint-calibration scalar $g$ (blue); the two agree without any shared parameter.}
  \label{fig:finger}
\end{figure}

\subsection{Quantifying the groundwater contribution: joint calibration}
Having detected the groundwater contribution model-free, we quantify it by jointly fitting the surface
and groundwater legs against the same data. The fit credits a statistically
significant groundwater term and closes the lower-mainstem residual: the fitted groundwater scalar is
$g=0.061$ (95\% CI $[0.022,\,0.100]$, excluding zero), significant against the surface-only ($g=0$)
model ($F_{1,5}=16.2$, $p=0.01$), and the mainstem log-RMSE falls from $0.15$ to $0.07$~dex
(Figure~\ref{fig:joint}; Table~\ref{tab:jointreach}). The improvement survives leave-one-out
cross-validation across the seven reaches (held-out RMSE $6.2\to4.5$~ng~L\textsuperscript{-1}), so it
is not a free parameter fitting in-sample. The groundwater contribution rises from $\sim$13\% of the
predicted concentration in the headwaters to $\sim$53\% at the lowest reach near the Wolverine
corridor, closing the sharp rise to $19$--$27$~ng~L\textsuperscript{-1} the surface model alone
under-predicts ($14.7$--$15.4$).

The joint fit is also the correct accounting. A naive sum of the two separately-calibrated legs
over-predicts the observed in-stream PFOS by about an order of magnitude because it double-counts: the
surface soil-loading factor was fitted to grab samples that already integrate the groundwater
contribution. Because in-stream concentration is linear in both the surface soil-loading and the
groundwater source strength, the combined prediction at mainstem reach $i$ is
$C_i = (L/L_0)\,C^{\mathrm{SW}}_i + g\,B_i$, where $C^{\mathrm{SW}}_i$ is the surface-water model's
in-stream PFOS at the reference loading $L_0=0.11$, $B_i$ is the cumulative groundwater-discharged
load routed to reach $i$ divided by its streamflow, and $(L,g)$ are fitted together by non-negative
least squares against the observed mainstem gradient in linear concentration space (the dex figures
quoted above are the derived diagnostic, not the fitted objective). (The linear superposition assumes in-stream
channel processes --- settling, resuspension, burial --- act linearly and identically on surface- and
groundwater-derived PFAS, a defensible first-order approximation given the low ambient sediment load.)
Fitting jointly drops the surface soil-loading from $L_0=0.11$ to $L=0.077$ --- the signature of
resolved double-counting, the groundwater leg now explaining part of the in-stream signal the
surface-only calibration had absorbed. The gain is not free: the joint fit slightly under-predicts the
headwaters and over-predicts the first lower-mainstem reach, consistent with a metric weighted toward
the high-concentration reaches, and the inference rests on only $n=7$ source-bearing reaches with $p=2$
free scalars. The reported significance ($p=0.01$) is moreover optimistic: consecutive mainstem reaches
share cumulative upstream load, so the seven residuals are spatially autocorrelated and the effective
sample size is below seven, inflating the nominal degrees of freedom and understating the true $p$. The
groundwater term is in practice carried by the two lowest reaches --- the only ones where the
groundwater fraction exceeds $50\%$ --- so the cross-validation ``survives'' largely because leaving one
of them out still leaves the other as an anchor, and the effective $n$ for the $g$ inference is closer to
two than to seven. (Because $g=0$ lies on the boundary of the non-negativity constraint, the standard
$F$ null is also only approximate, an error that is conservative here.) The result remains the most
defensible reading of these data, but it is observation-\emph{consistent} support, not a strong
significance claim.

Two cautions bound the interpretation. First, $g$ is a single lumped scalar rescaling the modeled
groundwater load; it absorbs source-magnitude error, interception, discharge geometry, and the
grid-controlled numerical dispersion (grid P\'eclet $\sim$25) together, and is not a clean physical
interception fraction. Second, the joint fit establishes that \emph{something} concentrated at the
lower mainstem is needed but on concentration alone cannot, by itself, distinguish groundwater
discharge from any co-located unmodeled term (an additional surface point source, a tributary input,
near-channel resuspension) loading the same reaches; the \emph{attribution} to groundwater is
established by the model-free multi-analyte fingerprinting above --- which a co-located surface source
with a different chemical signature cannot mimic --- and is corroborated by the aquifer-side validation. A correct
georeferencing of the groundwater discharge to the channel network was essential to expose this
signal: an earlier model-grid-local mapping mislocated the discharge to a headwater tributary and
produced a spurious null ($g\approx0$) that the spatially-explicit joint test caught but the
aquifer-only validation could not.

\begin{figure}[ht]
  \centering
  \includegraphics[width=0.85\textwidth]{joint_calibration.png}
  \caption{Joint surface-water + groundwater calibration on the Rogue mainstem (upstream to
  downstream). The joint fit (blue) tracks the observed rise (black) where the surface-water model
  alone (orange) under-predicts; the fitted groundwater contribution (green bars) is small in the
  headwaters and large at the lower mainstem near the Wolverine source, closing the lower-mainstem
  residual ($g=0.061$; joint $0.07$ vs surface-only $0.15$~dex).}
  \label{fig:joint}
\end{figure}

\paragraph{Predictive uncertainty.} Propagating the joint-calibration parameter uncertainty (the fitted
$(L,g)$ and their covariance) through the linear in-stream model gives a $95\%$ prediction interval at
each source-bearing mainstem reach (Figure~\ref{fig:uncertainty}). The observed PFOS lies within the
predictive band at every reach (log-space $R=0.95$, log-RMSE $0.07$~dex, percent bias $-1\%$), and the
band widens toward the lower mainstem where the groundwater contribution --- and its uncertainty --- is
largest. This is the predictive uncertainty conditional on the calibrated model structure; the broader
parametric uncertainty across the full coupled-model parameter set is screened by the global
sensitivity analysis below.

\begin{figure}[ht]
  \centering
  \includegraphics[width=\textwidth]{fig_uncertainty_obs_vs_sim.png}
  \caption{Predictive uncertainty of the joint surface-water + groundwater in-stream PFOS model on the
  seven source-bearing mainstem reaches. (a) Observed versus simulated PFOS (log axes) with the $1{:}1$
  line, the one-order-of-magnitude envelope, and the $95\%$ prediction interval (horizontal bars)
  propagated from the joint-calibration parameter uncertainty; $R=0.95$, log-RMSE $0.07$~dex, percent
  bias $-1\%$. (b) The same reaches upstream to downstream, with the $95\%$ prediction band, the
  simulated profile, and the observations; the band widens toward the groundwater-dominated lower
  mainstem.}
  \label{fig:uncertainty}
\end{figure}

\subsection{Global sensitivity: what controls the coupled prediction}
A calibration tunes parameters but does not reveal which of the coupled model's many parameters
actually govern its predictions, nor whether the surface and groundwater compartments interact. We
therefore screened the full coupled parameter space with the Morris elementary-effects method: $18$
parameters spanning the surface-water PFAS engine (soil loading, the in-stream organic-carbon
partition, the Langmuir air--water terms, the percolation partition), the groundwater flow model
(recharge and five conductance/conductivity multipliers), and the groundwater transport model
(Freundlich $K_f$ and $n$, the dispersivities, source magnitude, and background), over $r=20$
trajectories ($380$ coupled evaluations) on a $64$-core cloud instance. Each evaluation runs the
entire chain --- SWAT\textsuperscript{+} for the land-derived in-stream load and MODFLOW~6 for the
groundwater flow-and-transport discharge --- and we report the mean absolute elementary effect
$\mu^*$ (importance) and its spread $\sigma$ (interaction and nonlinearity) for four quantities of
interest (Table~\ref{tab:morris}; Supplementary Figure~S16). Thirteen percent of the samples ---
groundwater parameter combinations with near-stalled flow that fail to converge or exceed the
transport time budget --- did not complete; we compute the elementary effects step-by-step so that
every other evaluation still informs the screening (a median of $17$ effects per parameter).

The screening reproduces, from global variance rather than a single fit, the spatial finding the
joint calibration reached. At the lower mainstem near the Wolverine source the in-stream PFOS is
controlled by the \emph{groundwater} pathway --- the prescribed source magnitude (the largest
$\mu^*$), the aquifer hydraulic conductivity, and the streambed conductance that meters the
discharge --- with the surface soil-loading a secondary control; at the mid mainstem the ranking
inverts and the surface soil-loading dominates, the groundwater terms falling to a few percent of
its influence. The groundwater plume concentration is governed by the transport parameters (source,
Freundlich $n$ and $K_f$) and the baseflow purely by the flow parameters (recharge, streambed and
aquifer conductivity) --- exactly the separation the model structure implies, and a check that the
coupling carries each signal to the right place. The interaction signal is itself diagnostic: for
the surface soil-loading $\sigma/\mu^*\approx0.5$, a near-linear and additive control, whereas for
the groundwater source and conductivity $\sigma/\mu^*\approx1$, the strong nonlinearity expected of
Freundlich sorption and a threshold-like discharge. This supports the linear superposition the joint
calibration assumed for the surface leg, and it flags the groundwater source and transport
parameters as the ones that warrant variance-based decomposition. A Sobol analysis on a random-forest
surrogate fit to the screening ensemble (Section~\ref{sec:disc}) confirms it: the lower-mainstem
prediction variance is dominated by streambed conductance, aquifer conductivity and source magnitude
(total-order $S_T=0.53$, $0.41$, $0.40$), while the mid-mainstem prediction is surface-controlled
(soil loading, $S_T=0.93$) --- the same spatial attribution from global variance.
